\section{实施示例}

\begin{example}[把晚饭后的学习重建为一种转变设计]
设想一位读者反复打算在晚饭后做技术阅读，但却总是进入刷手机的状态，并持续整个晚上。

\textbf{步骤 1：诊断失败。}
\begin{itemize}[nosep]
  \item 当前状态：饭后浏览手机，具有快速新奇性和低努力特征。
  \item 目标状态：二十分钟带批注的阅读，并写下一个问题。
  \item 主导障碍项：当前奖赏过高、第一步存在较高歧义，以及饭后收拾结束后的窗口很窄。
\end{itemize}

\textbf{步骤 2：识别证据车道。}
\begin{itemize}[nosep]
  \item 强证据车道支持：状态转变研究使我们有理由把“晚饭到夜间”这一边界视为有意义的机制变化机会。
  \item 中等证据车道支持：亚稳态与吸引子式推理支持这样一种预期，即反复使用同样的线索与同样的书桌配置，会让未来的进入更加可靠。
  \item 设计推断：下面的具体配置属于应用设计，而不是直接引用结论。
\end{itemize}

\textbf{步骤 3：重设计这次转变。}
\begin{enumerate}[nosep]
  \item \textbf{窗口捕获：} 把干预点定义为洗完碗后紧接着的两分钟。
  \item \textbf{环境再加权：} 在开饭前就把手机移出房间，这样旧的夜间吸引子会更难重建。
  \item \textbf{状态预载：} 预先把书翻开，并贴上一张便签，写明要回答的一个问题。
  \item \textbf{最小可行进入：} 读一页，并写下一条页边批注。
  \item \textbf{历史塑形：} 连续七次使用同一把椅子、同一盏灯，以及同一个笔记本起始动作。
  \item \textbf{验证：} 记录第一条页边批注是否在晚上 8:15 前出现。
\end{enumerate}

\textbf{为什么这比泛泛建议更好。}
这一设计并不要求读者“感觉自己更认真”。它改变的是转变的局部几何结构。晚饭后的边界被设计成一个决策窗口；目标状态在进入之前已经被部分搭建；旧状态变得更难恢复；而重复的历史则被用来提高下一天晚上从更有利局部配置开始的概率。
\end{example}